\clearpage\pagestyle{otMain}\setlength{\otBlocEspacement}{2.5mm}\fontsize{11}{15}\selectfont\corrigesChapitre{9}{Les décimaux jusqu’aux millièmes}
\begin{blocLivre}{cor-9-1}\begin{corrige}{1}{Lire les rangs}\demarche{Compter les rangs après la virgule.}
Centièmes : $\res7$ ; millièmes : $\res5$. Décomposition : $\res{4+\frac3{10}+\frac7{100}+\frac5{1\,000}}$.
\end{corrige}\end{blocLivre}
\begin{blocLivre}{cor-9-2}\begin{corrige}{2}{Des fractions aux virgules}\demarche{Extraire la partie entière puis coder.}
$\res{2{,}375}$ ; $\res{0{,}208}$ ; $\res{0{,}045}$. Les zéros tiennent les rangs vides.
\end{corrige}\end{blocLivre}
\begin{blocLivre}{cor-9-3}\begin{corrige}{3}{Des virgules aux fractions}\demarche{Lire chaque chiffre avec son rang.}
$6{,}04=\res{6+\frac4{100}}$ : aucun dixième. $0{,}208=\res{\frac2{10}+\frac8{1\,000}}$ : aucun centième.
\end{corrige}\end{blocLivre}
\begin{blocLivre}{cor-9-4}\begin{corrige}{4}{Mètres et millimètres}\demarche{Placer les millimètres au rang des millièmes.}
1 mm $=0{,}001$ m : $\res{3{,}045}$ m ; $\res{2{,}007}$ m ; $\res{0{,}605}$ m.
\end{corrige}\end{blocLivre}
\begin{blocLivre}{cor-9-5}\begin{corrige}{5}{Le piège des longueurs d’écriture}\demarche{Compléter avec des zéros pour aligner les rangs.}
$6{,}400>6{,}385$ : $\res{6{,}4>6{,}385}$. Et $0{,}720>0{,}702$ : $\res{0{,}72>0{,}702}$. Plus de chiffres ne signifie pas plus grand.
\end{corrige}\end{blocLivre}
\begin{blocLivre}{cor-9-6}\begin{corrige}{6}{Ranger}\demarche{Penser chaque nombre en millièmes.}
En millièmes : $3\,052$ ; $3\,055$ ; $3\,205$ ; $3\,250$ ; $3\,500$. Donc $\res{3{,}052<3{,}055<3{,}205<3{,}25<3{,}5}$.
\end{corrige}\end{blocLivre}
\begin{blocLivre}{cor-9-7}\begin{corrige}{7}{Encadrer au dixième}\demarche{Geler les rangs avant celui demandé.}
$\res{8{,}2<8{,}236<8{,}3}$ et $\res{8{,}23<8{,}236<8{,}24}$.
\end{corrige}\end{blocLivre}
\begin{blocLivre}{cor-9-8}\begin{corrige}{8}{Intercaler}\demarche{Ajouter un rang décimal.}
En millièmes : $\res{5{,}275}$ convient. Entre $0{,}999$ et 1, penser en dix-millièmes : $\res{0{,}9995}$ convient. Entre deux décimaux distincts, il y a toujours de la place.
\end{corrige}\end{blocLivre}
\begin{blocLivre}{cor-9-9}\begin{corrige}{9}{Arrondir}\demarche{Regarder le chiffre des dixièmes.}
$7{,}493\to\res7$ (4 dixièmes) ; $7{,}51\to\res8$ ; $12{,}5\to\res{13}$, à mi-chemin on choisit l’entier supérieur.
\end{corrige}\end{blocLivre}
\begin{blocLivre}{cor-9-10}\begin{corrige}{10}{Sur la demi-droite}\demarche{Lire le pas avant les positions.}
$\res{A(4{,}21)}$ et $\res{B(4{,}27)}$. C se place au milieu du pas entre $4{,}23$ et $4{,}24$ ; en vert au corrigé.
\label{fig-cor-9-10}\figureGeo[corrige]{dec-10}
\end{corrige}\end{blocLivre}
\begin{blocLivre}{cor-9-11}\begin{corrige}{11}{Les performances}\demarche{Comparer en millièmes de mètre.}
$3{,}250>3{,}209>3{,}205$ : $\res{\text{Lou, puis Mia, puis Aya}}$. Écart de tête : $3{,}25-3{,}209=\res{0{,}041}$ m, soit 41 millimètres.
\end{corrige}\end{blocLivre}
\begin{blocLivre}{cor-9-12}\begin{corrige}{12}{Le même nombre ?}\demarche{Comparer rang par rang, pas à l’œil.}
$0{,}5=0{,}500=\frac5{10}$ forment un groupe ; $0{,}050=\frac{50}{1\,000}$ l’autre : dix fois plus petit. Les zéros finaux ne changent rien, les zéros intercalés si.
\end{corrige}\end{blocLivre}
\begin{blocLivre}{cor-9-13}\begin{corrige}{13}{La bonne précision}\demarche{Convertir en millimètres pour compter.}
$0{,}125$ m $=\res{125}$ mm. Dans $1\,000$ mm : $125\times8=1\,000$, soit $\res8$ pièces exactement.
\end{corrige}\end{blocLivre}
\begin{blocLivre}{cor-9-14}\begin{corrige}{14}{L’arrondi trompeur}\demarche{Chercher deux nombres proches d’un même entier.}
$\res{\text{Oui}}$ : $7{,}499$ et $6{,}502$ s’arrondissent tous deux à 7, pourtant $7{,}499>6{,}502$. L’arrondi estime, mais la comparaison exige les valeurs exactes.
\end{corrige}\end{blocLivre}
\begin{blocLivre}{cor-9-15}\begin{corrige}{15}{Zéro virgule neuf cent quatre-vingt-dix-neuf}\demarche{Compléter les écritures au même rang.}
$0{,}999>0{,}99$ est $\res{\text{vrai}}$ : 999 millièmes contre 990. Mais $0{,}999<1$ : il manque un millième. « Plus de chiffres » n’est jamais l’argument : seuls les rangs comptent.
\end{corrige}\end{blocLivre}
